\section{The proof}
\label{sec:proof}

The proof has the classical shape of a simulation argument
\cite{lynch1995simulations,abadi1991refinement}: a relation between the
states of the two systems, invariants of each side, a step simulation in
each direction, and then the lifting of those steps to finite and infinite
runs. Figure~\ref{fig:architecture} shows the parts. This section describes
them in order and says where each one is in the Lean development.

\begin{figure}[t]
\centering
\begin{tikzpicture}[
  >={Stealth[length=4.5pt]},
  box/.style={draw,rounded corners=2pt,align=center,font=\scriptsize,inner sep=4pt,minimum height=0.8cm,text width=3.0cm},
  thm/.style={box,very thick,fill=black!6,text width=3.7cm},
  x=1cm,y=1cm]
\node[box] (model) at (0,5.4) {Models and translation\\ \texttt{Dex/}, \texttt{L/}, \texttt{Translation/}};
\node[box] (linv) at (-4.3,3.9) {\Lang{} invariant\\ \texttt{L.Invariant.step} (WP1)};
\node[box] (rel) at (0,3.9) {Simulation state\\ \texttt{IsSimulationState}};
\node[box] (dinv) at (4.3,3.9) {Dex invariants\\ \texttt{ConfigInvariant}, \ldots};
\node[box] (ltod) at (-2.6,2.3) {\Lang{} to Dex steps\\ (WP2, WP3)};
\node[box] (dtol) at (2.6,2.3) {Dex to \Lang{} steps\\ (WP4)};
\node[box] (fin) at (-4.3,0.7) {Finite runs\\ \texttt{FiniteTraces.lean}};
\node[box] (blocks) at (0,0.7) {Block runs, silent weight\\ (L1, L2)};
\node[box] (inf) at (4.3,0.7) {Infinite runs, fairness\\ (L3 to L6)};
\node[thm] (safety) at (-4.3,-0.9) {\texttt{safety\_of\_isWellFormed}};
\node[thm] (live) at (4.3,-0.9) {\texttt{liveness\_of\_isWellFormed}};
\node[box,text width=2.6cm] (disc) at (0,-0.9) {Lock discipline,\\ witnesses};
\draw[->] (model) -- (linv);
\draw[->] (model) -- (rel);
\draw[->] (model) -- (dinv);
\draw[->] (linv) -- (ltod);
\draw[->] (rel) -- (ltod);
\draw[->] (rel) -- (dtol);
\draw[->] (dinv) -- (dtol);
\draw[->] (dinv) -- (ltod);
\draw[->] (linv) -- (dtol);
\draw[->] (ltod) -- node[font=\scriptsize,above] {determinism} (dtol);
\draw[->] (ltod) -- (fin);
\draw[->] (dtol) -- (fin);
\draw[->] (ltod) -- (blocks);
\draw[->] (dtol) -- (blocks);
\draw[->] (blocks) -- (inf);
\draw[->] (dtol) -- (inf);
\draw[->] (fin) -- (safety);
\draw[->] (inf) -- (live);
\draw[->,densely dashed] (disc) -- (safety);
\draw[->,densely dashed] (disc) -- (live);
\end{tikzpicture}
\caption{The structure of the proof. WP1 to WP4 are the work packages of
safety, L1 to L6 those of liveness (\S\ref{sec:engineering}). The dashed
arrows are non-vacuity (\S\ref{sec:nonvacuity}): the hypothesis holds for
concrete programs, and the trace sets are inhabited.}
\label{fig:architecture}
\end{figure}

\subsection{The simulation state}
\label{sec:relation}

Every step lemma takes one hypothesis, \lean{IsSimulationState}
(\file{Simulation/Relation.lean}), with five parts:

\begin{Verbatim}
structure IsSimulationState (program : Program) (lConfig : L.Config)
    (dexConfig : Dex.Config signature) : Prop where
  isRelated : Related program lConfig dexConfig
  lInvariant : L.Invariant program lConfig
  dexInvariant : ConfigInvariant program dexConfig
  tasksDistinct : TasksAreDistinct dexConfig
  timersDistinct : TimersAreDistinct dexConfig
\end{Verbatim}

\paragraph{\lean{Related}} holds only the clauses that mention both sides.
The clocks agree, and the last stamped \Lang{} action carries the timestamp
of the last Dex commit: this is the timing model at work. The statuses
correspond (not started, running, closed). While the program runs
(\lean{RunningRelated}):
\begin{itemize}
\item \emph{Data.} Shared variables are Attributes; each condition
  variable's queue is its waiter-list Attribute; each external channel's
  queue is its Channel's pending messages.
\item \emph{Tokens.} A mutex's token Channel holds one token exactly while
  the mutex is unlocked; $\mathit{join}/t$ holds one message exactly while
  $t$ has ended and was not joined; $\mathit{cv}/w$ holds the wakes of a
  waiter that has not parked yet. A pending cancel request means the flag
  is set; a set flag means a pending request, or a thread that has observed
  it, or an ended thread that no live thread can still fork.
\item \emph{Timers.} The armed Dex timers are the armed \Lang{} timers,
  with the same number and deadline, and the timer counter is the counter
  of the \lean{timed} Step type.
\item \emph{Threads and Steps.} Every thread that has not halted is carried
  by exactly one active Step, and every active Step carries one thread
  (\lean{StepCarriesThread}): the Step has the type, input, and options the
  translation gives that thread. The thread's phase fixes the Step's phase,
  invocation, and task (\lean{PhaseMatchesStep}), as in
  Figure~\ref{fig:lifecycle}. In particular a Step is in phase
  $\mathit{WT}$ exactly while its thread is blocked, so a publication wakes
  a Step exactly when it wakes the thread.
\item \emph{Order.} Two blocked threads are ordered by their blocking
  stamps as their Steps' waits are ordered by $\mathit{sws}$. This carries
  FIFO from one side to the other.
\end{itemize}
After the close (\lean{ClosedRelated}), the late emits of \Lang{} are, up to
permutation, the effects that the remaining stale tasks have not emitted
yet; no Step has an invocation, and no timer is armed.

\paragraph{The invariants.}
\lean{L.Invariant} (\file{L/Invariant.lean}) holds the facts about \Lang{}
alone: thread IDs are unique and the next children of a thread are fresh;
the next mutex a thread creates does not exist yet; every handle to a
\emph{created} mutex, or to a condition variable of one, that a thread
can reach names a mutex that \kw{newMutex} has created (this is where
unforgeable handles are used);
blocking stamps and timer numbers are distinct; each thread's phase
fits its code; a thread that began a statement that requires a mutex holds
it until the statement completes. The invariant does not check a static
index against the program's number of static mutexes:
\kw{staticMutex}$(i)$ with $i$ out of range names a mutex that never
exists, and a \kw{lock} of it blocks forever on both sides. On the Dex
side,
\lean{ConfigInvariant} wraps a run invariant (\lean{Dex.Run.Invariant}:
unique Step IDs, counted Step numbers, injective $\mathit{sws}$, only the
wait shapes the translation uses, and no pending message on a Channel that
a waiting Step names) and says that a task holds the response its handler
computes; \lean{TasksAreDistinct} and \lean{TimersAreDistinct} say
that task IDs and armed timer IDs are unique. Each invariant is proved to
hold initially and to be kept by every transition
(\lean{L.Invariant.step}, \lean{ConfigInvariant.step},
\lean{TasksAreDistinct.step}, \lean{TimersAreDistinct.step}).

\subsection{From \Lang{} to Dex: named action sequences}
\label{sec:ltodex}

The \Lang{}-to-Dex direction does not say ``there exist Dex actions''. It
names them. \lean{simulatingActions} is a function of the two
configurations and the \Lang{} label:

\begin{itemize}
\item a tick, a request, and a batch of \kw{fire} actions are one Dex action
  each (\lean{tick}, \lean{request}, \lean{timerFire});
\item a thread that begins an action is the \rulename{STEP-START} of its
  Step;
\item $\ev{out}(a)$ is the \lean{workerEffect} of the thread's task, after
  the \lean{workerRun} that runs the handler if it has not run;
\item a thread that completes an action is the \lean{stepComplete} of its
  task, after the \lean{workerRun} and \lean{workerReturn} that have not
  happened yet;
\item after the close, the delivery or the drop of a late emit is the
  \lean{workerEffect} or the \lean{workerAbandon} of the first task that
  still holds the effect.
\end{itemize}
So one \Lang{} label is one to three Dex actions. The step lemma, proved by
cases on the label (\lean{l_step_simulatedByActions},
\file{Simulation/LToDex.lean}), says that these actions are enabled, have
the same projected events, and end in a related configuration:

\begin{Verbatim}
def IsSimulatedByActions (program : Program) (lConfig : L.Config)
    (dexConfig : Dex.Config signature) (label : L.Label) (lEvents : List Event)
    (lNext : L.Config) : Prop :=
  ∃ dexEvents dexNext,
    (Dex.system ⟦program⟧).exec dexConfig (simulatingActions program lConfig dexConfig label) =
      some (dexEvents, dexNext) ∧
    dexEvents.filterMap project = lEvents ∧
    Related program lNext dexNext
\end{Verbatim}

The bulk of the proof is here. Registering a wait
(\file{Simulation/Register.lean}) and completing an action
(\file{Simulation/Completion.lean} and \file{Simulation/Completion/}) must
re-establish every clause of \lean{Related} after a commit that may write
an Attribute, publish tokens, wake the earliest waiter, create Steps, or
close the Flow.

\subsection{From Dex to \Lang{}: at most one label, by determinism}
\label{sec:dextol}

In the other direction, every Dex step of an admissible run is simulated by
at most one \Lang{} label (\lean{dex_step_simulated},
\file{Simulation/DexToL.lean}):

\begin{Verbatim}
def IsSimulatedByL (program : Program) (lConfig : L.Config)
    (dexEvents : List (Dex.Obs signature)) (dexNext : Dex.Config signature) : Prop :=
  ∃ labels lNext,
    labels.length ≤ 1 ∧
    (L.system program).exec lConfig labels = some (dexEvents.filterMap project, lNext) ∧
    Related program lNext dexNext
\end{Verbatim}

Running a handler, returning it, and dropping a task that holds no effect
have no \Lang{} label: the list is empty and the relation is kept. For the
actions that \Lang{} sees, the proof does not repeat the case analysis of
the relation. It reuses the other direction:
\begin{enumerate}
\item Show that the corresponding \Lang{} label is enabled.
\item Apply the \Lang{}-to-Dex lemma. It names the simulating Dex actions.
\item Show that those actions are the given action alone. For a Worker
  effect or a completion, the given action is enabled only when the Worker
  steps before it have already happened, so the named list has no prefix.
\item \lean{Dex.step} is a function. The given step is therefore the
  simulating step, and its result is the related configuration of step 2.
\end{enumerate}
This is why the \Lang{}-to-Dex statement names its actions: an existential
statement could not be matched against a given Dex step. Compiler
verification uses determinism for a similar economy: with a deterministic
target, preservation proved in one direction suffices
\cite{leroy2009compcert}.

Step~1 is the only place where well-formedness is used. At
\rulename{STEP-START} of a Step whose statement requires a mutex, the
\Lang{} thread can begin only if it holds the mutex; the hypothesis that
\lean{BreaksMutexRequirement} does not hold says that it does. Consequently the
two \Lang{}-to-Dex inclusions, \lean{finiteTrace_l_to_dex} and
\lean{fairMaximalTrace_l_to_dex}, hold for every program; only the
Dex-to-\Lang{} inclusions need the scope.

\subsection{Safety}

Safety follows by induction on a finite run
(\file{FiniteTraces.lean}). The initial configurations are a simulation
state (\lean{related_initial} and the initial cases of the invariants). Each
step of one side is matched by the step simulation, and the invariants are
carried along by their own step lemmas. The Dex-to-\Lang{} induction also
carries reachability of the \Lang{} configuration, to apply
well-formedness, and admissibility of each Dex action, which follows from
the admissible prefix. The two inclusions are \lean{finiteTrace_dex_to_l}
and \lean{finiteTrace_l_to_dex}.

\subsection{Liveness}
\label{sec:liveness-proof}

Each inclusion of liveness builds an infinite run of one side from an
infinite run of the other, and then carries over the input word, the trace,
the marker, and the fairness assumptions.

\paragraph{Runs in blocks.}
The two systems do not step in lockstep. One \Lang{} label is one to three
Dex actions, and a Dex Worker step is no \Lang{} label. A run of one side
therefore gives the other side a \emph{run in blocks}
(\lean{BlockRun}, \file{Core/BlockRun.lean}): one finite, possibly empty,
list of labels per step of the first (Figure~\ref{fig:blocks}). A block run
with infinitely many non-empty blocks \emph{flattens} into an ordinary
infinite run (\lean{BlockRun.IsFlattenedTo}), and the flattening keeps the
trace, the attempted requests, and whether the progress steps are finitely
many. This part mentions neither Dex nor \Lang{}.

\begin{figure}[t]
\centering
\begin{tikzpicture}[
  >={Stealth[length=4pt]},
  cell/.style={draw,minimum height=0.55cm,minimum width=0.62cm,font=\scriptsize,inner sep=1pt},
  sil/.style={cell,fill=black!12},
  lab/.style={font=\scriptsize},
  blk/.style={draw,densely dotted,rounded corners=2pt,inner sep=2.5pt},
  x=1cm,y=1cm]
% (a) L run -> Dex blocks
\node[lab,anchor=west] at (-1.6,2.9) {(a) A Dex run built from an \Lang{} run: blocks of \texttt{fairActions}};
\node[lab,anchor=east] at (-0.1,2.2) {\Lang{}};
\node[cell] (a0) at (0.9,2.2) {begin};
\node[cell] (a1) at (3.4,2.2) {$\mathit{out}$};
\node[cell] (a2) at (6.2,2.2) {complete};
\node[cell] (a3) at (8.6,2.2) {tick};
\node[cell] (a4) at (10.4,2.2) {exit};
\node[lab,anchor=east] at (-0.1,1.2) {Dex};
\node[cell] (p0) at (0.9,1.2) {start};
\node[sil] (p1) at (3.05,1.2) {run};
\node[cell] (p2) at (3.75,1.2) {fx};
\node[sil] (p3) at (5.85,1.2) {ret};
\node[cell] (p4) at (6.55,1.2) {cmpl};
\node[cell] (p5) at (8.6,1.2) {tick};
\node[sil] (p6) at (9.7,1.2) {ret};
\node[cell] (p7) at (10.4,1.2) {cmpl};
\node[sil] (p8) at (11.1,1.2) {drop};
\node[blk,fit=(p0)] (g0) {};
\node[blk,fit=(p1)(p2)] (g1) {};
\node[blk,fit=(p3)(p4)] (g2) {};
\node[blk,fit=(p5)] (g3) {};
\node[blk,fit=(p6)(p7)(p8)] (g4) {};
\draw[->] (a0) -- (g0);
\draw[->] (a1) -- (g1);
\draw[->] (a2) -- (g2);
\draw[->] (a3) -- (g3);
\draw[->] (a4) -- (g4);
% (b) Dex run -> L blocks
\node[lab,anchor=west] at (-1.6,0.2) {(b) An \Lang{} run built from a Dex run: blocks of length 0 or 1};
\node[lab,anchor=east] at (-0.1,-0.5) {Dex};
\node[cell] (q0) at (0.9,-0.5) {start};
\node[sil] (q1) at (2.3,-0.5) {run};
\node[cell] (q2) at (3.7,-0.5) {fx};
\node[cell] (q3) at (5.1,-0.5) {tick};
\node[sil] (q4) at (6.5,-0.5) {ret};
\node[cell] (q5) at (7.9,-0.5) {cmpl};
\node[sil] (q6) at (9.3,-0.5) {stale};
\node[lab,anchor=east] at (-0.1,-1.5) {\Lang{}};
\node[cell] (r0) at (0.9,-1.5) {begin};
\node[lab] (r1) at (2.3,-1.5) {$[\,]$};
\node[cell] (r2) at (3.7,-1.5) {$\mathit{out}$};
\node[cell] (r3) at (5.1,-1.5) {tick};
\node[lab] (r4) at (6.5,-1.5) {$[\,]$};
\node[cell] (r5) at (7.9,-1.5) {complete};
\node[lab] (r6) at (9.3,-1.5) {$[\,]$};
\foreach \a/\b in {q0/r0,q1/r1,q2/r2,q3/r3,q4/r4,q5/r5,q6/r6} \draw[->] (\a) -- (\b);
\end{tikzpicture}
\caption{Runs in blocks. Dotted outlines are blocks. Shaded Dex actions are
of a \emph{silent kind}:
\Lang{} has no label for them. (a) Each \Lang{} label becomes its
simulating actions, followed, after the close, by the drops of the canceled
tasks that hold no effect. (b) Each Dex action becomes at most one \Lang{}
label; silent actions become empty blocks.}
\label{fig:blocks}
\end{figure}

\paragraph{Dex to \Lang{}: the silent weight.}
Given a fair Dex run, \lean{dex_step_simulated} gives an \Lang{} run in
blocks of length 0 or 1 (\lean{LFollowsDex}). The danger is an infinite
tail of empty blocks: the Dex run would have infinitely many progress steps
(\mdiv) while the \Lang{} run has finitely many. The proof rules it out
with a measure. The actions that can have no \Lang{} label are of four
kinds: \lean{workerRun}, \lean{workerReturn}, \lean{workerAbandon}, and
\lean{stepCompleteStale}. A task can take at most three such steps (run,
return, drop), and the \emph{silent weight} of a configuration counts what
is left:

\begin{Verbatim}
def Dex.Task.silentWeight {σ : Dex.Signature} (task : Dex.Task σ) : Nat :=
  match task.phase with
  | .started => 3
  | .running _ _ => 2
  | .returned _ => 1

def Dex.Config.silentWeight {σ : Dex.Signature} (config : Dex.Config σ) : Nat :=
  (config.tasks.map Dex.Task.silentWeight).sum
\end{Verbatim}

Every action of a silent kind strictly lowers the weight
(\lean{silentWeight_lt_of_isSilentKind}), a tick keeps it, and only a step
that \Lang{} sees can raise it. So infinitely many progress steps of the
Flow include infinitely many that \Lang{} sees, and the \Lang{} labels of a
Dex step are progress steps exactly when the action is one. The markers
then agree. F1 gives infinitely many ticks, hence infinitely many non-empty
blocks, so the flattening exists (\lean{exists_lRun_of_dexRun}).

\paragraph{\Lang{} to Dex: cleaning up after the close.}
Given a fair \Lang{} run, the simulating actions give a Dex run in blocks
(\lean{DexFollowsL}). One more step is needed for fairness. When the
program ends, the Flow close cancels every active Step, and the tasks of
those Steps stay in $\Theta$. A stale task that still holds an effect
corresponds to a late emit, and \Lang{}'s own fairness delivers or drops
it. A stale task that holds no effect has no \Lang{} step at all, yet F4
requires it to leave $\Theta$. So the Dex run does not use
\lean{simulatingActions} alone. It uses \lean{fairActions}: the simulating
actions, followed by the \rulename{WORKER-ABANDON} of every canceled task
that holds no effect (\lean{cleanupActions}). Between two blocks no such
task and no \emph{Returned} task is left.

\paragraph{Fairness and markers carry over.}
For each direction, one lemma per assumption shows that the built run is
fair when the given run is. The correspondences are those of
Table~\ref{tab:fairness}: a thread whose next action has not begun is a
pending invocation (F3); a thread that issued an emit or began an action is
a task, with an effect exactly when it issued an emit (F4, F5); a due timer
is a due timer (F2); a late emit is an effect of a stale task (F4). The
markers carry over because closed corresponds to closed in \lean{Related}
and because flattening keeps whether the progress steps are finitely many.
The two inclusions are \lean{fairMaximalTrace_dex_to_l} and
\lean{fairMaximalTrace_l_to_dex} (\file{Liveness.lean}).

\subsection{Non-vacuity}
\label{sec:nonvacuity}

A theorem with a hypothesis is only as good as the set of objects that
satisfy the hypothesis. \lean{Program.IsWellFormed} quantifies over every
reachable configuration, so a concrete trace does not establish it, and
without a witness the theorems could hold for no program at all.

\paragraph{A static lock discipline.}
\file{L/LockDiscipline.lean} defines a decidable, syntactic, sufficient
condition, \lean{Program.hasStaticLockDiscipline}: a must-hold lock-set
analysis over the statically known mutexes.
\kw{lock}$(\kw{staticMutex}(i))$ adds $i$ to the set. \kw{unlock},
\kw{wait}, \kw{signal}, and \kw{broadcast} need a statically known operand
whose mutex is in the set. The two branches of an \kw{if} or a \kw{select}
must agree. A loop body must give back the set it started with. A called or
forked procedure body, a \kw{try} body, and a handler are checked from the
empty set. \file{L/LockDisciplineSound.lean} proves soundness:

\begin{Verbatim}
theorem Program.isWellFormed_of_hasStaticLockDiscipline (program : Program)
    (isDisciplined : program.hasStaticLockDiscipline = true) : program.IsWellFormed
\end{Verbatim}

The proof rests on one fact of \Lang{}: only the owner of a mutex releases
it, so the steps of other threads never shrink the set of mutexes a thread
holds.

\paragraph{Well-formed programs exist.}
\file{Examples/WellFormed.lean} proves \lean{IsWellFormed} for eight
programs, each by evaluating the discipline in the Lean kernel. Together
they use every construct of \Lang{}: a two-thread mutex program; a program
with \kw{sleep}, a timed \kw{select}, \kw{recv}, \kw{cancel}, \kw{fork},
and \kw{join}; a \kw{wait}/\kw{signal} program; a
\kw{wait}-in-a-loop/\kw{broadcast} program with shared variables; a
recursive program; a program with \kw{try}, a \kw{throw} across a call, and
an uncaught exception; three threads racing in \kw{select}; and a program
that creates and locks a mutex and then fails.
\file{Examples/WellFormedSafety.lean} states the safety theorem at three of
them with no hypothesis left, for example:

\begin{Verbatim}
theorem mutexProgram_safety (trace : List Event) :
    IsProjectedDexFiniteTrace Mutex.program trace ↔ IsLFiniteTrace Mutex.program trace :=
  safety_of_isWellFormed Mutex.program mutexProgram_isWellFormed trace
\end{Verbatim}

The trace sets are not empty either. \file{Examples/Mutex.lean} shows, by
running explicit schedules in the kernel, that both complete traces of the
two-thread mutex program are finite traces of the program and projected
finite traces of its Flow (\lean{bothTraces_onBothSides}).
\file{Examples/TimedSelect.lean} and \file{Examples/TimerOrder.lean} do the
same for timers, \kw{select}, cancellation, and two timers with equal
deadlines.

\paragraph{Fair maximal traces exist.}
The trace sets of liveness are defined by the existence of a fair run.
Without a fair run of some program, \lean{liveness_of_isWellFormed} would
compare two empty sets. \file{Liveness/RestingRuns.lean} builds fair
maximal runs of \Lang{}. A configuration is \emph{at rest}
(\lean{L.Config.isAtRest}) when every thread is halted or blocked without a
timer and no late emit is left; then none of F2 to F5 speaks about it. A
run that reaches a configuration at rest by a finite schedule, and then
takes only steps that keep it at rest while time diverges, is fair
(\lean{IsRestingTail.isLFairMaximalTrace}). Two such tails are built: clock
ticks alone, and clock ticks alternating with inputs that are refused after
the close.

\file{Examples/FairRuns.lean} gives the witnesses of
Table~\ref{tab:fairruns}. Each finite schedule is checked by evaluation in
the kernel. Each \Lang{} trace is then also a projected fair maximal trace
of the translated Flow, by instantiating \lean{liveness_of_isWellFormed} at
a well-formed program and an admissible word. So both sides of the liveness
equivalence are inhabited for each marker.

\begin{table}[t]
\centering
\small
\begin{tabularx}{\textwidth}{@{}P{1.2cm}P{3.5cm}P{2.7cm}L@{}}
\toprule
Marker & Program & Input word & \Lang{} side; Dex side \\
\midrule
\mterm & the two-thread mutex program & the start &
  \lean{mutexProgram_terminates}; \lean{mutexProgram_terminates_dex} \\
\mdead & \kw{emit}$(1)$; \kw{recv}$(5)$ & the start &
  \lean{waitingProgram_deadlocks}; \lean{waitingProgram_deadlocks_dex} \\
\mdead & \kw{lock}$(m_0)$; \kw{emit}$(1)$; \kw{lock}$(m_0)$ & the start &
  \lean{selfDeadlockProgram_deadlocks};
  \lean{selfDeadlockProgram_deadlocks_dex} \\
\mdiv & the two-thread mutex program & the start, then one input forever &
  \lean{mutexProgram_diverges}; \lean{mutexProgram_diverges_dex} \\
\bottomrule
\end{tabularx}
\caption{Kernel-checked fair maximal traces
(\file{Examples/FairRuns.lean}). The Dex-side theorems follow from the
\Lang{}-side ones by the liveness theorem.}
\label{tab:fairruns}
\end{table}

One limit of these witnesses must be stated precisely. The \mdiv{} witness
has an \emph{infinite input word}: the program exits, and every later input
is refused, forever. Those refusals are the infinitely many progress steps.
It is not a program that diverges on its own. No witness is checked for a
program that runs forever on a finite word, such as a loop that keeps
emitting: its configurations never come to rest, so its run is not a finite
schedule followed by a regular tail. There is also no general lemma that
every reachable configuration of a well-formed program has a fair maximal
continuation.
